% LSPV-1 REPORT::INFONOGO1::2026-10-12T17:11:00Z
\documentclass[11pt]{article}
\usepackage{fontspec}
\usepackage{geometry}
\usepackage{amsmath,amssymb,mathtools}
\usepackage{booktabs,tabularx}
\usepackage{enumitem}
\usepackage[hidelinks]{hyperref}
\geometry{margin=1in}
\setlist{nosep}
\newcommand{\OPEN}{\textsf{OPEN}}
\newcommand{\EXACT}{\textsf{PROVED/EXACT}}
\newcommand{\COMP}{\textsf{COMPUTATIONAL}}
\newcommand{\RET}{\textsf{FALSIFIED/RETIRED}}
\newcommand{\same}{\textnormal{\textsc{same\_realization}}}
\title{INFONOGO1: Compatibility Mutual Information Is Not an Unrestricted Clean-Core Lower Bound}
\author{LSPV-1 / Universal-core lower bound}
\date{13 October 2026}

\begin{document}
\sloppy
\maketitle

\section*{Status}
\textbf{P versus NP remains open.}
\textsf{DETERMINISTIC-TARGET-INFORMATION-TRILEMMA1} is \EXACT{} as an
elementary Shannon-information theorem, pending independent review.  Its
finite identity audit is \COMP{}.
\textsf{MUTUAL-INFORMATION-COMPATIBILITY-LB1} is \RET{}.
\textsf{DIRECT-CLEANCORE-TIME-LB1} remains \OPEN{}.

\section*{Corpus and literature preflight}
This run read the current repository protocol, root \path{main.tex}, actual
\path{main4.tex} and \path{main5.tex}, the newest MAIN synthesis
\textsf{TRACECONJ1}, and the newest verified LSPV-2, LSPV-3, and LSPV-4
reports \textsf{SERIALTRACE1}, \textsf{CAUSALCANCEL1}, and
\textsf{MODELBRIDGE-0028}.  The immediately preceding LSPV-4 canonical
artifact verifies.  Primary literature was checked first for Shannon
information, data processing, SAT time--space lower bounds, relativization,
algebrization, and Natural Proofs.

\section*{Explicit clean core and exact lower-bound target}
Let $\mathcal K$ be the fixed NP-complete signed Monotone
$3$SAT-$(2,2)$ language under a fixed self-delimiting binary encoding:
clauses have three distinct variables and one sign, duplicate clauses are
excluded, and every variable occurs twice positively and twice negatively.
For $x=\langle F,\alpha\rangle$ of encoded length $N$, define
\[
 H(F,\alpha)=\{\beta:\alpha\cup\beta\models F\},\qquad
 h(F,\alpha)={\bf 1}[H(F,\alpha)\ne\varnothing].
\]
The direct target, with no unnecessary extension to unrestricted formula
families, is
\[
 \forall A\in\operatorname{Dec}(\mathcal K)\;
 \forall d\ge1\;\forall N_0\;\exists x\in\operatorname{Dom}(\mathcal K):
 |x|\ge N_0\ \land\ T_A(x)>|x|^d.
\]
Proving this statement would give $\mathcal K\notin\mathsf P$ and hence
$\mathsf P\ne\mathsf{NP}$.  It is \OPEN{}.  This run tests whether a
distributional compatibility-information obligation can be a weaker route to
that target.

\section*{Candidate under attack}
Let $X$ be distributed on one fixed-length slice of the clean core, let
$H=g(X)$ be a selected compatibility object or bit, let $A(X)$ be the one-bit
decision, and let $T=t(X)$ be a retained view of one deterministic realized
execution.  A proposed bridge is:
\begin{quote}
Every correct decider must generate substantial mutual information between
$T$ and $H$, beyond the answer bit; that information forces time, space, or
communication cost.
\end{quote}
This is attractive because mutual information ignores coordinate names and
many representation changes.  It fails because $H$ is a deterministic
function of the input source, and Shannon information contains no
computational provenance.

\section*{Theorem: DETERMINISTIC-TARGET-INFORMATION-TRILEMMA1}
For every finite random variable $X$, deterministic target $H=g(X)$,
deterministic quotient $A=a(X)$, and deterministic trace view $T=t(X)$:
\begin{enumerate}
\item source-conditioned target information is identically zero,
\[
 I(H;T\mid X)=I(H;T\mid X,A)=0;
\]
\item conditional data processing gives
\[
 I(H;T\mid A)\le I(X;T\mid A),
\]
so excluding all residual source information also excludes residual
compatibility information; and
\item a source copy or any invertible source transport attains the maximum,
\[
 T=X\quad\Longrightarrow\quad
 I(H;T\mid A)=H(H\mid A).
\]
The maximum may therefore be produced by $O(N)$ copying without an
$H$-specific operation.
\end{enumerate}

\paragraph{Proof.}
Because $H=g(X)$, $H(H\mid X)=H(H\mid X,A)=0$, proving (1).  Conditional on
$A$, the variables satisfy $H\leftarrow X\rightarrow T$, so the conditional
data-processing inequality proves (2).  If $T$ is invertible to $X$, then
$H(H\mid T,A)=0$, and
\[
 I(H;T\mid A)=H(H\mid A)-H(H\mid T,A)=H(H\mid A),
\]
proving (3). \hfill$\square$

\section*{Exact clean-core witness}
Use eight equal-length satisfiable 3-CNFs that force the eight assignments
$u\in\{0,1\}^3$, represented with fixed-width literal blocks.  On their
uniform distribution the decision $A$ is constant YES while $H$ is the unique
forced assignment, so $H(H\mid A)=3$ bits.  Both a raw source copy and the
invertible prefix-parity trace satisfy
\[
 I(H;T\mid A)=3\text{ bits},
\]
although the added trace mechanism is only source transport.  A parity
projection carries one bit.  Every trace satisfies
$I(H;T\mid X)=0$.  Thus neither large unconditional information nor
information beyond the decision bit certifies compatibility computation.

\section*{Why the natural repairs fail}
\paragraph{Normalize by source information.}
The ratio $I(H;T\mid A)/I(X;T\mid A)$ can equal one for a trace that merely
retains one target-correlated source bit, and is undefined when both terms
vanish.  It supplies no runtime lower bound.

\paragraph{Subtract copied information.}
Since $H$ is itself a function of $X$, the joint distribution does not
canonically label a bit as copied or semantically computed.  COPY, RECOMP,
INTERP, delay, and alternate encodings can have the same law for
$(X,H,A,T)$ with different dataflow and cost.

\paragraph{Condition on a source descriptor.}
If the descriptor omits information relevant to $H$, source leakage survives.
If it retains all $H$-relevant source information, the conditional target
information vanishes.  Proving that an adequate descriptor is hard is itself
an algorithmic lower bound, not a Shannon identity.

\paragraph{Convert bits to time or space.}
Entropy is not time, peak workspace, communication, circuit size, or proof
length.  No model-matched unconditional theorem converts the displayed bits
into a superpolynomial cost for arbitrary uniform deterministic deciders.

\section*{Bounded reproducible audit}
The script
\path{experiments/review/deterministic_target_information_trilemma_audit.py}
enumerates uniform sources $\{0,1\}^n$ for $1\le n\le8$, six targets, three
decision quotients, and seven trace projections.  It verifies:
\begin{itemize}
\item $2{,}016$ conditional-zero identities;
\item $1{,}008$ conditional data-processing inequalities;
\item $390$ source-free implications;
\item $432$ maximum-information COPY/prefix/invertible-encoding identities;
\item $360$ hostile-control assertions across nine families, five wrappers,
and four encodings; and
\item the five exact forced-assignment values $3,3,3,1,0$ bits.
\end{itemize}
All $4{,}206$ enumerated assertions and the five forced-assignment values
passed.  Results are stored at
\path{results/raw/review/2026-10-13_deterministic_target_information_trilemma.json}
with SHA-256
\texttt{ed832efc996de130e728cc499629d4026a680811b3ebaa428f98ad0b7bb83bf5}.
The audit is finite and no asymptotic runtime is inferred from it.

\section*{Hostile controls and costs}
The exact identities cover 2SAT, Horn, affine/Tseitin, bounded-treewidth,
independent-product, padding, redundant-clause, alternate-encoding, parity,
and signed-clean-core controls.  Tractable semantic shortcuts confirm rather
than weaken the no-go: implication graphs, Horn propagation, Gaussian
elimination, and bounded-width dynamic programming can compute selected
targets cheaply, while source copying transports them without computing them.

All parsing, input acquisition, copy/prefix generation, observation,
extraction, reconstruction, and output costs are separate.  A full source
copy costs $O(N)$ time and $O(N)$ explicit storage; a serial prefix trace has
linear generation cost.  Neither charge becomes superpolynomial.  Compact
alternate representations and the one-bit decision output remain admissible.

\section*{Faithful relation to main4/main5 and peers}
The actual \path{main4.tex} is the ARC achievable-region manuscript and the
actual \path{main5.tex} is the WRAPPER construction manuscript.  Neither
claims an entropy-to-runtime theorem for all clean-core deciders.  Their
conditional controllers still require all-input acquisition/closure and a
model-matched landing.

\textsf{SERIALTRACE1} supplies the invertible prefix source transport.
\textsf{CAUSALCANCEL1} blocks causal liveness as a provenance repair.
\textsf{MODELBRIDGE-0028} blocks data-free schedule plus answer.
\textsf{TRACECONJ1} shows those guards can coexist in one polynomial wrapper.
The present theorem retires Shannon target information as the next generic
repair.

\section*{Barrier and counterfactual audit}
The theorem relativizes: an oracle execution can still copy its source, and
$H=g(X)$ remains deterministic.  It uses no arithmetization and supplies no
nonalgebrizing hinge.  No large constructive circuit property is given, so it
neither invokes nor evades Natural Proofs.  Circuit, branching-program,
communication, proof-system, and SAT time--space lower bounds remain confined
to their declared endpoint models.

Assume the clean core has a uniform polynomial-time decider $A$.  A wrapper
can copy $X$, run $A$, and attain maximum $I(H;T\mid A)$ in polynomial time.
Selected compatibility queries can also be answered by separately charged
restriction/reduction calls to $A$.  Both are consistent with the assumption.
Declaring either cost superpolynomial would smuggle the desired conclusion
into a premise.

\section*{Transfer ledger and surviving arrow}
\begin{center}
\begin{tabularx}{\textwidth}{p{0.33\textwidth}p{0.20\textwidth}X}
\toprule
Claim & Status & Scope\\
\midrule
Deterministic-target information trilemma & \EXACT{} & Finite random
variables; review pending\\
Finite identity/hostile-control audit & \COMP{} & $n\le8$; no asymptotics\\
Mutual information as paid compatibility work & \RET{} & COPY and
source-conditioning countermodels\\
Direct clean-core superpolynomial time lower bound & \OPEN{} & All uniform
deterministic deciders\\
Model-matched unconditional contradiction & \OPEN{} & No landing supplied\\
P versus NP & \OPEN{} & No proof\\
\bottomrule
\end{tabularx}
\end{center}

The missing implication is not transfer to broader SAT inputs.  A genuine
all-algorithm lower bound for this fixed NP-complete core would suffice.  The
missing mathematical content is instead either:
\begin{enumerate}
\item a machine-level $H$-specific resource invariant that cannot be created
by source transport or erased by wrappers, together with an unconditional
superpolynomial lower bound for that exact same realization; or
\item a direct diagonal, combinatorial, or complexity-theoretic lower bound
for every uniform deterministic decider of $\mathcal K$.
\end{enumerate}
Both are \OPEN{}.

\section*{Peer handoff}
LSPV-2 should test time-bounded or computational-information repairs against
streaming copies, lazy reconstruction, and universal-machine dependence.
LSPV-3 should attack directed-information and causal-provenance candidates
with cancellation and state recoding.  LSPV-4 should reject any external
information extractor whose work is not charged to the same execution.
LSPV-5 should add the entropy trilemma to the retired universal-transfer
ledger while preserving the direct NP-complete-core lower-bound route.

\section*{Primary references checked}
\begin{itemize}
\item C. E. Shannon, ``A mathematical theory of communication,'' Bell System
Technical Journal 27 (1948), 379--423 and 623--656,
\url{https://people.math.harvard.edu/~ctm/home/text/others/shannon/entropy/entropy.pdf}.
\item T. Baker, J. Gill, and R. Solovay, ``Relativizations of the $P=?NP$
question,'' SIAM Journal on Computing 4(4), 1975, 431--442,
\url{https://doi.org/10.1137/0204037}.
\item S. Aaronson and A. Wigderson, ``Algebrization: a new barrier in
complexity theory,'' ACM Transactions on Computation Theory 1(1), 2009,
\url{https://doi.org/10.1145/1490270.1490272}.
\item A. A. Razborov and S. Rudich, ``Natural proofs,'' Journal of Computer
and System Sciences 55(1), 1997, 24--35,
\url{https://doi.org/10.1006/jcss.1997.1494}.
\item L. Fortnow, R. Lipton, D. van Melkebeek, and A. Viglas,
``Time-space lower bounds for satisfiability,'' Journal of the ACM 52(6),
2005, 835--865, \url{https://doi.org/10.1145/1101821.1101822}.
\end{itemize}

% END_OF_VERIFIED_REPORT::LSPV-1::INFONOGO1
\end{document}
